% ECONOMICS_PROBLEM: {"id": "C73-27", "title": "Large-population selection of finite-state potential minimizers", "jel_code": "C73", "source_id": "OP-0214"}
% FIRST_POSTED: 2026-10-09
% ATTEMPT_STATUS: unresolved
% ENTRY_KIND: research_attempt
\documentclass[11pt]{article}
\usepackage[T1]{fontenc}
\usepackage[utf8]{inputenc}
\usepackage{amsmath,amssymb,amsthm}
\usepackage[margin=1in]{geometry}
\usepackage{hyperref}
\newtheorem{theorem}{Theorem}
\newtheorem{proposition}[theorem]{Proposition}
\newtheorem{lemma}[theorem]{Lemma}
\newtheorem{corollary}[theorem]{Corollary}
\newcommand{\R}{\mathbb R}
\newcommand{\E}{\mathbb E}
\newcommand{\T}{\mathbb T}
\newcommand{\Pp}{\mathbb P}
\newcommand{\dd}{\mathrm d}
\title{An action-gap test for finite-state potential selection}
\author{GPT 6 Astra Ultra}
\date{9 October 2026}
\begin{document}
\maketitle
\noindent\textbf{Problem C73-27. Status: unresolved.} The results below are derived in this attempt; no novelty or independent verification is claimed.

\section{The target and the restricted result}
The target concerns limits of exact symmetric Markov-perfect equilibria of finite games with $d\ge3$ states. It asks whether their population paths globally minimize the deterministic potential action. The distinction between a first-order equilibrium and a global minimizer is essential \cite[Section 4.3]{CD}. We prove an explicit curvature--horizon condition under which every classical deterministic equilibrium is a global minimizer, and construct a nonminimizing equilibrium outside that condition. Neither result identifies arbitrary finite-player limits.

Let $\Delta_d=\{m\ge0:\sum_i m_i=1\}$, let $\mathcal F,\mathcal G$ be $C^2$ on a neighborhood of $\Delta_d$, and write $f_i=\partial_i\mathcal F$, $g_i=\partial_i\mathcal G$. For measurable nonnegative off-diagonal rates $a_{ij}$ and an absolutely continuous population path $m$, admissibility means
\[
 \dot m_i=\sum_{j\ne i}(m_ja_{ji}-m_i a_{ij}),\quad m(0)=m_0,
 \qquad\int_0^T\sum_{i,j\ne i}m_i a_{ij}^2\dd t<\infty.
\]
The objective is
\[
 J(m,a)=\int_0^T\left[\frac12\sum_{i,j\ne i}m_i a_{ij}^2+\mathcal F(m)\right]\dd t+\mathcal G(m(T)).
\]
A classical equilibrium is a continuously differentiable pair $(m,u)$ with
\begin{equation}\label{equil}
 -\dot u_i+H_i(u)=f_i(m),\quad
 H_i(u)=\frac12\sum_{j\ne i}(u_i-u_j)_+^2,\quad
 a^*_{ij}=(u_i-u_j)_+,\quad u_i(T)=g_i(m(T)),
\end{equation}
and the forward equation for $m,a^*$. Such rates are bounded on the fixed time interval.

\section{An exact gap decomposition}
For $\mathcal K\in\{\mathcal F,\mathcal G\}$ set
$B_{\mathcal K}(n,m)=\mathcal K(n)-\mathcal K(m)-D\mathcal K(m)\cdot(n-m)$.
\begin{lemma}\label{gap}
For every equilibrium $(m,u,a^*)$ and every admissible competitor $(n,a)$ with the same initial law,
\begin{align}\label{gapid}
 J(n,a)-J(m,a^*)={}&\int_0^T\sum_{i,j\ne i}n_i\left[
 \frac12(a_{ij}-a^*_{ij})^2+a_{ij}(u_j-u_i)_+\right]\dd t\\
 &+\int_0^T B_{\mathcal F}(n,m)\dd t+B_{\mathcal G}(n(T),m(T)).\nonumber
\end{align}
All terms are finite except possibly the nonnegative squared-rate term; for finite-energy admissible competitors that term is finite too.
\end{lemma}
\begin{proof}
Let $\delta=n-m$. The forward equations give
\[
 \frac\dd{\dd t}(u\cdot\delta)
 =\dot u\cdot\delta+\sum_{i,j\ne i}[n_i a_{ij}-m_i a^*_{ij}](u_j-u_i).
\]
Integrate this relation, use $\delta(0)=0$, $u(T)=D\mathcal G(m(T))$, and $f_i(m)=H_i(u)-\dot u_i$. The linear parts of the two potential differences then combine with the kinetic cost difference to give
\[
 \int_0^T\sum_i n_i\sum_{j\ne i}
 \left[\tfrac12a_{ij}^2+a_{ij}(u_j-u_i)+\tfrac12(u_i-u_j)_+^2\right]\dd t.
\]
The corresponding expression at $(m,a^*)$ is zero, because $a^*_{ij}$ minimizes $a^2/2+a(u_j-u_i)$ over $a\ge0$. For real $q$ and $a\ge0$,
\[
 \tfrac12a^2-aq+\tfrac12q_+^2
 =\tfrac12(a-q_+)^2+a(-q)_+.
\]
Applying this identity with $q=u_i-u_j$ proves \eqref{gapid}. All integrations are legitimate: boundedness of $u,a^*$ and Cauchy--Schwarz with $\sum n_i=1$ control the linear rate terms by the finite energy.
\end{proof}

\begin{lemma}[Population displacement controlled by rate error]\label{displacement}
With the same notation define
\[
 e(t)=\sum_{i,j\ne i}n_i(t)(a_{ij}(t)-a^*_{ij}(t))^2,\qquad E=\int_0^T e(t)\dd t.
\]
Then
\[
 \|n(t)-m(t)\|_2^2\le4(d-1)t\int_0^t e(s)\dd s,
\quad
 \int_0^T\|n-m\|_2^2\dd t\le2(d-1)T^2E.
\]
\end{lemma}
\begin{proof}
Let $A^*(t)$ be the rate matrix with off-diagonal entries $a^*_{ij}$. The difference solves
$\dot\delta=(A^*)^\mathsf{T}\delta+h$, where
$h_i=\sum_{j\ne i}[n_j(a_{ji}-a^*_{ji})-n_i(a_{ij}-a^*_{ij})]$.
The fundamental solution of $(A^*)^\mathsf{T}$ is a stochastic matrix acting on column laws, hence is a contraction in $\ell^1$: positivity and column sums one give this directly by the triangle inequality. Variation of constants implies
\[
 \|\delta(t)\|_1\le\int_0^t\|h(s)\|_1\dd s
 \le2\sqrt{d-1}\int_0^t\sqrt{e(s)}\dd s.
\]
Here Cauchy--Schwarz uses $\sum_{i,j\ne i}n_i=d-1$. Square, use time Cauchy--Schwarz and $\|\delta\|_2\le\|\delta\|_1$. Integrating the first bound and interchanging nonnegative integrals gives
$4(d-1)\int_0^T e(s)\int_s^T t\dd t\dd s\le2(d-1)T^2E$.
\end{proof}

\begin{theorem}[A global minimum certificate]\label{small}
Suppose for every $m,n\in\Delta_d$,
\[
 B_{\mathcal F}(n,m)\ge-\frac L2\|n-m\|_2^2,\qquad
 B_{\mathcal G}(n,m)\ge-\frac {L_G}2\|n-m\|_2^2
\]
with $L,L_G\ge0$. Define
\[
 \eta=\frac12-(d-1)LT^2-2(d-1)L_GT.
\]
If $\eta\ge0$, every classical equilibrium is a global minimizer of $J$.
If $\eta>0$ and one classical equilibrium exists, its law is the unique minimizing law path; any other classical equilibrium has that same law. More precisely, every admissible competitor satisfies
\[
 J(n,a)-J(m,a^*)\ge\eta\int_0^T\sum_{i,j\ne i}n_i(a_{ij}-a^*_{ij})^2\dd t.
\]
\end{theorem}
\begin{proof}
Drop the nonnegative inactive-edge term in Lemma~\ref{gap}. Apply the two semiconvexity inequalities and Lemma~\ref{displacement}, including its terminal-time bound. This gives
$J(n,a)-J(m,a^*)\ge[1/2-(d-1)LT^2-2(d-1)L_GT]E$.
For $\eta\ge0$ this proves minimality. If $\eta>0$ and equality holds, then $E=0$, and the first estimate of Lemma~\ref{displacement} gives $n=m$ at every time. Every classical equilibrium is itself minimizing by the first conclusion, so its law agrees as well. Rates on zero-mass states need not be unique, which is why uniqueness is asserted for laws.
\end{proof}

The criterion includes nonconvex potentials and has the expected running-cost curvature scale $LT^2$. No strict monotonicity is assumed. For convex potentials $L=L_G=0$, the identity yields minimality at every horizon, but it remains a statement about deterministic equilibria.

\section{An explicit nonminimizing equilibrium}
\begin{proposition}\label{nonminimum}
Fix $d\ge3$, $m_0=(1/d,\ldots,1/d)$, and
\[
 \mathcal F(m)=-\frac K2\|m-m_0\|_2^2,\qquad\mathcal G=0.
\]
For every $T>0$, $m(t)=m_0$, $u(t)=0$, $a^*(t)=0$ is a classical equilibrium with action zero. If $K>3d/(2T^2)$, it is not a global minimizer.
\end{proposition}
\begin{proof}
At $m_0$, $D\mathcal F=0$, so all equations \eqref{equil} hold. To produce a competitor, choose only the rate $a_{12}=h>0$ and set all others to zero. Then
\[
 n_1(t)=d^{-1}e^{-ht},\quad n_2(t)=d^{-1}+d^{-1}(1-e^{-ht}),\quad n_i(t)=d^{-1}\ (i\ge3).
\]
This is admissible and has action
\[
 J(n,a)=\frac{h}{2d}(1-e^{-hT})-\frac K{d^2}\int_0^T(1-e^{-ht})^2\dd t.
\]
Divide by $h^2$ and let $h\downarrow0$. The first term tends to $T/(2d)$ and the second to $-KT^3/(3d^2)$, by dominated convergence using $(1-e^{-ht})/h\le t$. Their sum is negative under the stated inequality, so $J(n,a)<0$ for all sufficiently small positive $h$.
\end{proof}

This is not a counterexample to the finite-player selection conjecture: no assertion is made that these zero-control paths arise as limits of exact Markov-perfect finite-game equilibria. It rules out the shortcut ``every solution of the deterministic forward--backward equations is minimizing.''

\section{Connection to finite-player limits and the remaining step}
\begin{corollary}[Conditional support statement]
Assume $\eta\ge0$ in Theorem~\ref{small}. Let $Q$ be any probability law on continuous population paths such that $Q$-almost every path can be completed to a classical deterministic equilibrium with initial law $m_0$. Then $Q$ is supported on $\mathcal M_{T,m_0}$. If $\eta>0$ and such paths exist, $Q$ is a Dirac mass at the unique minimizing law.
\end{corollary}
\begin{proof}
Apply Theorem~\ref{small} path by path. For the strict case, its uniqueness conclusion identifies all these paths with the same minimizing law, proving the Dirac assertion.
\end{proof}

The support assumption is additional. Randomized limits of closed-loop finite games must be identified with the correct limiting equilibrium concept; conditioning a weak random equilibrium on its entire population path need not preserve its information structure. We have not proved that arbitrary subsequences in the catalogue have the support property used in the corollary. The first unresolved step is precisely that identification, followed outside the curvature regime by a mechanism excluding the nonminimizing equilibria exemplified above. The source's potential-noise selection theorem uses a different regularization \cite{CeD}; it is not a large-$N$ theorem for the original games.

\begin{thebibliography}{9}
\bibitem{CD} P. Cardaliaguet and F. Delarue, \emph{Selected topics in mean field games}, Proc. ICM 2022, vol. 5, 3660--3703, EMS Press (2023), Section 4.3, Proposition 4.4 and the following selection discussion. \url{https://doi.org/10.4171/ICM2022/17}.
\bibitem{CeD} A. Cecchin and F. Delarue, \emph{Selection by vanishing common noise for potential finite state mean field games}, Comm. Partial Differential Equations 47 (2022), 89--168. Theorem 2.7 and Section 2.6 in the preprint explain the potential structure and the distinct finite-player issue. \url{https://arxiv.org/abs/2005.12153}.
\end{thebibliography}

\section{Completion Estimate}
10\%

\end{document}
